\chapter{Finite Hamiltonian jets and relative quantization}
\label{ch:vol2-hamiltonian-jets}

A Hamiltonian determines an infinitesimal motion of a symplectic space.
To study a boundary fixed at a chosen point, we keep the Hamiltonians
whose vector fields vanish there. Their finite jets give Lie algebras
with an explicit quantization. The resulting associative algebras have
relative Hochschild complexes, so they also give finite models of the
coherent actions considered in Chapter~\ref{ch:vol2-boundary-actions}.

The lowest nonzero jet is already instructive. Quadratic Hamiltonians form
$\mathfrak{sl}_2$; cubic Hamiltonians supply its four-dimensional
representation; the bracket of two cubics first survives at quartic
order. We shall construct these algebras, their Poisson cochains and
their finite Batalin--Vilkovisky actions before comparing the
quantized algebra with its Hochschild cochains.

\section{Hamiltonians fixing a point}

Work over $\C$ on the formal plane with coordinates $z_1,z_2$ and
symplectic form $\omega=dz_1\wedge dz_2$. Set
$\mathfrak m=(z_1,z_2)\subset\C\llbracket z_1,z_2\rrbracket$.
Our Poisson convention is
\begin{equation}\label{eq:vol2-hamiltonian-bracket}
 \{f,g\}=\partial_{z_1}f\,\partial_{z_2}g
             -\partial_{z_2}f\,\partial_{z_1}g.
\end{equation}
Thus the vector field $X_f(g)=\{f,g\}$ vanishes at the origin
exactly when the linear part of $f$ vanishes. Constants give the zero
vector field. Choosing the representative with zero constant term
identifies the pointed Hamiltonian Lie algebra with
$\h=\mathfrak m^2$.

For the quadratic polynomials
\[
 e=\tfrac12z_1^2,\qquad f=-\tfrac12z_2^2,\qquad h=-z_1z_2,
\]
direct differentiation gives
\begin{equation}\label{eq:vol2-quadratic-brackets}
 \{h,e\}=2e,\qquad \{h,f\}=-2f,\qquad \{e,f\}=h.
\end{equation}
They are the Chevalley generators of $\mathfrak{sl}_2$.

Write $R_r$ for the homogeneous polynomials of degree $r$.
Each derivative lowers degree by one, so
$\{R_r,R_s\}\subset R_{r+s-2}$.
The filtration
\[
 F^p\h=\mathfrak m^{p+2}\quad(p\geq0)
\]
is therefore additive:
$\{F^p\h,F^q\h\}\subset F^{p+q}\h$.
It is separated and complete in the formal power-series topology.
For $d\geq2$, define
\begin{equation}\label{eq:vol2-finite-jet}
 \h_{\leq d}=\h/\mathfrak m^{d+1}
             =\bigoplus_{r=2}^{d}R_r.
\end{equation}
The equality on the right specifies the polynomial representatives;
their bracket is projected to degrees at most $d$.

\begin{proposition}[Finite Hamiltonian jets]\label{prop:vol2-finite-jets}
The bracket makes $\h_{\leq d}$ a Lie algebra of dimension
\[
 \dim\h_{\leq d}=\frac{d(d+3)}2-2.
\]
Its quadratic subalgebra is $\mathfrak{sl}_2$, and
$R_r$ is its irreducible representation of highest weight $r$.
\end{proposition}

\begin{proof}
The ideal $\mathfrak m^{d+1}$ is stable under bracketing with
$\mathfrak m^2$. Hence Jacobi descends from the Poisson bracket.
In coordinates, this bracket is induced by the constant bivector
$\partial_{z_1}\wedge\partial_{z_2}$; commuting partial derivatives
cancel the Jacobiator. Its monomial formula is
\begin{equation}\label{eq:vol2-monomial-bracket}
 \{z_1^az_2^b,z_1^cz_2^s\}
 =(as-bc)z_1^{a+c-1}z_2^{b+s-1},
\end{equation}
where a term with zero coefficient is zero, including when a displayed
exponent would be negative. For three exponent vectors
$\alpha,\beta,\gamma$ and $\epsilon=(1,1)$, the Jacobi coefficient is
\[
 \det(\alpha,\beta)\det(\alpha+\beta-\epsilon,\gamma)
 +\text{cyclic permutations}=0.
\]
The terms quadratic in the pairwise determinants cancel. The remaining
terms vanish by the two-dimensional identity
\[
 \det(\alpha,\beta)\gamma+\det(\beta,\gamma)\alpha
 +\det(\gamma,\alpha)\beta=0.
\]
The dimension is $\sum_{r=2}^d(r+1)=d(d+3)/2-2$.

Put $u_j=z_1^{r-j}z_2^j$ for $0\leq j\leq r$. The quadratic action is
\begin{equation}\label{eq:vol2-weight-action}
 \{h,u_j\}=(r-2j)u_j,\qquad
 \{e,u_j\}=ju_{j-1},\qquad
 \{f,u_j\}=(r-j)u_{j+1}.
\end{equation}
Terms outside the index range are zero. A nonzero invariant subspace
contains a weight vector because the eigenvalues of $h$ are distinct.
Repeated application of $e$ reaches $u_0$, and repeated application of
$f$ then reaches every $u_j$. This proves irreducibility.
\end{proof}

The pointed condition is needed for these quotients. The affine algebra
\[
 \h^{\mathrm{aff}}=\C\llbracket z_1,z_2\rrbracket/\C
\]
also contains linear Hamiltonians, which generate translations. In it,
\[
 \{z_1,z_2^{d+1}\}=(d+1)z_2^d.
\]
Consequently $\mathfrak m^{d+1}$ is no longer an ideal. A moving
boundary or a centre-of-mass coordinate requires an affine construction
in addition to the pointed jet quotient.

\section{The first nonlinear brackets}

For $d=2$ the quotient is $\mathfrak{sl}_2$ itself. For $d=3$, set
$v_j=z_1^{3-j}z_2^j$ for $0\leq j\leq3$.
Equation~\eqref{eq:vol2-weight-action} gives the quadratic action on
their span.

\begin{proposition}[The cubic jet]\label{prop:vol2-cubic-jet}
The degree-three quotient is the semidirect product
\[
 \h_{\leq3}=\mathfrak{sl}_2\ltimes R_3,
 \qquad [R_3,R_3]=0.
\]
\end{proposition}

\begin{proof}
A bracket of two cubics has degree four and vanishes in this quotient.
The brackets with quadratics are exactly the action in
\eqref{eq:vol2-weight-action}.
\end{proof}

At $d=4$, put $w_k=z_1^{4-k}z_2^k$ for $0\leq k\leq4$.
The cubic bracket becomes
\begin{equation}\label{eq:vol2-cubic-quartic}
 \{v_i,v_j\}=3(j-i)w_{i+j-1}.
\end{equation}
When the index is outside $0,\ldots,4$, the coefficient is zero.
In particular,
\[
 \{v_0,v_3\}=9w_2,\qquad \{v_1,v_2\}=3w_2.
\]
The first equality differentiates $z_1^3$ and $z_2^3$.
The second is
$4z_1^2z_2^2-z_1^2z_2^2=3z_1^2z_2^2$.
The two cubic pairs reach the same quartic monomial with coefficients
in the ratio $3:1$.

\begin{proposition}[The quartic jet]\label{prop:vol2-quartic-jet}
On the vector space $R_2\oplus R_3\oplus R_4$, all brackets are
determined by \eqref{eq:vol2-quadratic-brackets},
\eqref{eq:vol2-weight-action} for $r=3,4$,
\eqref{eq:vol2-cubic-quartic}, and
\[
 [R_3,R_4]=0,\qquad [R_4,R_4]=0.
\]
This Lie algebra is $\h_{\leq4}$. Its quadratic and cubic elements
generate it; its ideal $R_3\oplus R_4$ is two-step nilpotent with
commutator subalgebra $R_4$.
\end{proposition}

\begin{proof}
The omitted brackets have degrees five and six, so they vanish in the
quotient. Every quartic is a cubic bracket: use the pairs
$(i,j)=(0,1),(0,2),(0,3),(1,3),(2,3)$ to obtain
$w_0,w_1,w_2,w_3,w_4$, respectively, with nonzero coefficients.
These formulas exhaust the brackets of a basis. Jacobi follows from
Proposition~\ref{prop:vol2-finite-jets}. They also give the stated
commutator and its centrality inside the ideal.
\end{proof}

The adjoint trace will measure the divergence of the finite cotangent
action. The homogeneous decomposition already determines this trace.

\begin{proposition}[Unimodularity]\label{prop:vol2-jet-unimodular}
For every $d\geq2$ and $x\in\h_{\leq d}$,
$\operatorname{Tr}(\operatorname{ad}_x)=0$.
\end{proposition}

\begin{proof}
A homogeneous Hamiltonian of degree greater than two strictly raises
polynomial degree. Its adjoint matrix is strictly triangular in the
decomposition \eqref{eq:vol2-finite-jet}. For a quadratic Hamiltonian,
use the basis $e,f,h$. On each $R_r$, the matrices of $e$ and $f$ have
zero diagonal, while the trace of $h$ is
$\sum_{j=0}^r(r-2j)=0$. Linearity proves the assertion.
\end{proof}

\section{Chevalley cochains of the linear Poisson structure}

Fix a finite-dimensional complex Lie algebra $\flie$, with basis
$e_a$ and bracket $[e_a,e_b]=f^c{}_{ab}e_c$; repeated indices are
summed. We shall apply the constructions to $\flie=\h_{\leq d}$.
Use cohomological grading, with $(W[k])^n=W^{n+k}$.
In particular, a vector of degree zero has degree one in $W[-1]$.

Let $\flie^\vee$ carry the coadjoint action
$(\operatorname{ad}_x^*\lambda)(y)=-\lambda([x,y])$.
Define the graded Lie algebra
\[
 \mathfrak q(\flie)=\flie\ltimes\flie^\vee[-1]
\]
with zero internal differential and abelian second summand.
Its elements in the second summand have degree one.
For a finite graded Lie algebra $\mathfrak q$, we use polynomial
Chevalley cochains
$C^*_{\CE}(\mathfrak q)=\Sym(\mathfrak q^\vee[-1])$
with differential dual to the bracket. Here $\Sym$ means the graded
symmetric algebra, without completion in its even generators.

\begin{proposition}[Chevalley--Poisson complex]\label{prop:vol2-ce-poisson}
There is a natural isomorphism of commutative differential graded
algebras
\[
 C^*_{\CE}(\mathfrak q(\flie))
 \cong C^*(\flie,\Sym\flie)
 =\Sym\flie\otimes\Lambda^*\flie^\vee.
\]
On the right the differential is the Chevalley differential for the
adjoint action. It is also the linear Poisson differential on
polynomial polyvector fields on $\flie^*$.
\end{proposition}

\begin{proof}
Dualizing and shifting the degree-one summand gives
$(\flie^\vee[-1])^\vee[-1]=\flie$ in degree zero.
The degree-zero summand gives exterior generators $c^a$ of degree one.
Let $x_a\in\Sym\flie$ denote the linear function
$x_a(\lambda)=\lambda(e_a)$. Dualizing the semidirect bracket gives
\begin{equation}\label{eq:vol2-ce-coordinate}
 d_{\CE}c^c=-\tfrac12 f^c{}_{ab}c^ac^b,
 \qquad d_{\CE}x_a=f^c{}_{ba}c^b x_c.
\end{equation}
These are precisely the bracket and action terms of the Chevalley
differential. On a polynomial $F$ the second formula reads
\[
 d_{\CE}F=c^a\{x_a,F\}_{\mathrm{lin}},
 \qquad \{x_a,x_b\}_{\mathrm{lin}}=f^c{}_{ab}x_c.
\]
Identify $c^a$ with the constant vector field $\partial_{x_a}$.
Use the Schouten bracket extending $[X,F]=X(F)$ by graded
skew-symmetry and the graded derivation rule.
For the bivector
$\pi=\tfrac12 f^c{}_{ab}x_c\partial_{x_a}\wedge\partial_{x_b}$,
the Schouten differential $[\pi,-]$ gives both formulas in
\eqref{eq:vol2-ce-coordinate}. Since both operators are derivations
of the exterior product, they agree on all polyvectors. Jacobi is
$[\pi,\pi]=0$, and hence $d_{\CE}^2=0$.
\end{proof}

\section{The finite BV action}

The preceding complex has degree-zero polynomial coordinates $x_a$.
The BV construction uses the different graded space
$T^*[-1](\flie[1])$. Its polynomial functions are
\[
 \mathcal B(\flie)=\Lambda(c^1,\ldots,c^m)
                  \otimes\C[b_1,\ldots,b_m],
 \qquad |c^a|=1,\quad |b_a|=-2,
\]
where $m=\dim\flie$. Thus the cotangent coordinates $b_a$ are
antifields of degree $-2$. Fix left derivatives in the odd variables
and set
\begin{equation}\label{eq:vol2-bv-conventions}
 \Delta=\sum_a\partial_{b_a}\partial_{c^a},\qquad
 \{F,G\}_{\mathrm{BV}}
 =(-1)^{|F|}\bigl(\Delta(FG)-\Delta(F)G
                    -(-1)^{|F|}F\Delta(G)\bigr).
\end{equation}
The bracket has degree one, with
$\{b_a,c^b\}_{\mathrm{BV}}=\delta_a^b$ and
$\{c^b,b_a\}_{\mathrm{BV}}=-\delta_a^b$.
Commutation of the even derivatives and anticommutation of the odd
ones give $\Delta^2=0$.
These Darboux coordinates specify a constant Berezinian up to a
nonzero scalar; its normalization does not change $\Delta$.
The quantum master equation for this measure is proved in
Theorem~\ref{thm:vol2-finite-qme}.

With these conventions the degree-zero cotangent action is
\begin{equation}\label{eq:vol2-bv-action}
 S=-\tfrac12 b_c f^c{}_{ab}c^ac^b.
\end{equation}
In particular, $\{S,c^a\}_{\mathrm{BV}}=d_{\CE}c^a$.

\begin{theorem}[Finite quantum master equation]\label{thm:vol2-finite-qme}
For every finite-dimensional Lie algebra $\flie$, the action
\eqref{eq:vol2-bv-action} satisfies
\[
 \tfrac12\{S,S\}_{\mathrm{BV}}=0,\qquad
 \Delta S=\sum_b\operatorname{Tr}(\operatorname{ad}_{e_b})c^b.
\]
Consequently, for $\flie=\h_{\leq d}$ and a formal parameter
$\hbarq$ of degree zero,
\[
 \tfrac12\{S,S\}_{\mathrm{BV}}+\hbarq\Delta S=0.
\]
This identity holds in $\mathcal B(\h_{\leq d})\llbracket\hbarq\rrbracket$.
\end{theorem}

\begin{proof}
Put $Q^a=-\tfrac12 f^a{}_{bc}c^bc^c$, so that $S=b_aQ^a$.
The coordinate bracket gives
\[
 \tfrac12\{S,S\}_{\mathrm{BV}}
 =\sum_{a,k}b_kQ^a\partial_{c^a}Q^k.
\]
The coefficient of $b_k$ is $d_{\CE}^2c^k$, which vanishes by
Jacobi. For the second identity, left differentiation gives
\[
 \Delta S=-\sum_{a,b}f^a{}_{ab}c^b
          =\sum_{a,b}f^a{}_{ba}c^b.
\]
The last coefficient is the trace of $\operatorname{ad}_{e_b}$.
Proposition~\ref{prop:vol2-jet-unimodular} makes it zero for every
finite Hamiltonian jet.
\end{proof}

The trace condition is substantive. In the two-dimensional Lie algebra
with basis $u,v$ and $[u,v]=v$, one has
\[
 S=-b_vc^uc^v,\qquad \Delta S=c^u.
\]
The classical master equation still holds, but this constant-measure
action fails the quantum master equation. For Hamiltonian jets,
unimodularity removes this finite-dimensional divergence.
A local field theory has further analytic and renormalization data;
the finite identity alone does not construct its quantum observables.

\section{Rees multiplication and interval operations}

Quantization replaces commuting linear functions on $\flie^*$ by
generators whose commutator is $\hbarq$ times the Lie bracket.
For $N\geq0$, let
\[
 R_N=\C[\hbarq]/(\hbarq^{N+1}),\qquad
 \flie_N=R_N\otimes_\C\flie,
 \qquad [x,y]_N=\hbarq[x,y].
\]
Here $N$ is the coefficient order, independent of the jet order $d$
or a Chan--Paton rank. Define the augmented $R_N$-algebra
\[
 A_N=U_{R_N}(\flie_N)
 =T_{R_N}(\flie_N)/(xy-yx-\hbarq[x,y]),
\]
whose augmentation sends every $x\in\flie$ to zero.
The Poincar\'e--Birkhoff--Witt (PBW) theorem has the following form
over this coefficient ring.

\begin{proposition}[PBW over the finite coefficient ring]\label{prop:vol2-rees-pbw}
Choose an ordered basis of $\flie$. Its ordered monomials form an
$R_N$-basis of $A_N$. If $S_0=\Sym_\C\flie$, symmetrization is a
filtered $R_N$-module isomorphism
\[
 \operatorname{sym}_N:S_0\otimes R_N\longrightarrow A_N,
 \qquad x_1\cdots x_r\longmapsto
 \frac1{r!}\sum_{\sigma\in S_r}x_{\sigma(1)}\cdots x_{\sigma(r)}.
\]
It intertwines the adjoint action through
\begin{equation}\label{eq:vol2-scaled-pbw-action}
 [x,\operatorname{sym}_N(F)]_{A_N}
 =\hbarq\operatorname{sym}_N(\operatorname{ad}_xF).
\end{equation}
\end{proposition}

\begin{proof}
Replace every inverted adjacent pair $yx$ by
$xy+\hbarq[y,x]$. In the first term the number of inversions
decreases; the second has shorter length. Induction on length and
then inversions proves termination. Disjoint replacements commute.
For an overlapping triple $z>y>x$, the difference between the two
reductions, after the shorter words have been ordered, is
$\hbarq^2$ times
\[
 [[z,y],x]+[[y,x],z]+[[x,z],y]=0.
\]
Induction therefore makes the ordered result independent of every
choice of replacement. This defines a normal-form map killing the
relations and equal to the identity on ordered monomials. Those
monomials are consequently a basis, even though $R_N$ has nilpotents.

Give generators PBW filtration degree one. Symmetrization induces
the identity on the associated graded polynomial algebra. Subtracting
the highest-degree symbol and descending in degree constructs its
inverse. This uses only finite degrees and the invertibility of
factorials in $R_N$. Finally, commuting $x$ past each factor of a
symmetrized word gives one term $\hbarq[x,x_i]$ in each position.
Their sum is \eqref{eq:vol2-scaled-pbw-action}.
\end{proof}

Transporting multiplication by $\operatorname{sym}_N$ gives the
associative PBW product on $S_0\otimes R_N$. Its reduction modulo
$\hbarq$ is the ordinary product. For $N\geq1$, the first commutator
is the linear Poisson bracket. Indeed, the coefficient of $\hbarq$ in a
commutator is a biderivation of the commutative reduction; its value
on $x,y\in\flie$ is $[x,y]$. These values determine the biderivation
on the polynomial algebra.

For disjoint intervals $I_1<\cdots<I_r$ contained in an interval $J$,
assign the operation
\[
 A_N^{\otimes_{R_N}r}\longrightarrow A_N,
 \qquad a_1\otimes\cdots\otimes a_r\longmapsto a_1\cdots a_r.
\]
Use the unit for an empty collection and the identity for an inclusion
of one interval. Associativity makes these operations compatible with
nested collections, and their independence of interval lengths makes
them locally constant. They supply the algebraic ordered interval
operations expected of a boundary quantization.
For a cotangent field theory on a half-plane, fix the cotangent field
at the boundary and leave the Lie algebra field free. Identifying the
resulting boundary observables with these interval operations requires
a map from that theory's quantum
observables, compatible with interval inclusions and fusion, and a
proof that the map is a quasi-isomorphism. PBW constructs the target
algebra, while this realization remains a separate comparison.

\section{The relative Hochschild comparison}

The coefficient ring affects the action complex. For the quadratic
jet, $[h,e]=2e$, whereas $[h,e]_{A_N}=2\hbarq e$.
Modulo $\hbarq$, the Hochschild differential of the zero-cochain $e$
vanishes because the algebra is commutative. The ordinary Chevalley
differential has $(d_{\CE}e)(h)=2e$. Thus it is the scaled differential
$\hbarq d_{\CE}$ that can compare with relative Hochschild cochains.

Write $\bar A_N$ for the augmentation ideal.
Use normalized relative Hochschild cochains
\[
 C^n_{\HH}(A_N/R_N,A_N)
 =\operatorname{Hom}_{R_N}(\bar A_N^{\tensor_{R_N}n},A_N)
\]
with differential
\[
\begin{split}
 (bF)(a_1,\ldots,a_{n+1})
 ={}&a_1F(a_2,\ldots,a_{n+1})\\
 &+\sum_{i=1}^n(-1)^iF(a_1,\ldots,a_i a_{i+1},\ldots,a_{n+1})\\
 &+(-1)^{n+1}F(a_1,\ldots,a_n)a_{n+1}.
\end{split}
\]
Let $d_{\CE}$ act on
$C^n(\flie,S_0)=\operatorname{Hom}_\C(\Lambda^n\flie,S_0)$
by the adjoint Chevalley differential.

\begin{theorem}[Relative Rees comparison]\label{thm:vol2-relative-rees}
Antisymmetrization defines a quasi-isomorphism of $R_N$-complexes
\[
 \operatorname{Alt}_N:
 C^\bullet_{\HH}(A_N/R_N,A_N)
 \longrightarrow
 \bigl(C^\bullet(\flie,S_0)\tensor R_N,\hbarq d_{\CE}\bigr),
\]
given by
\[
 (\operatorname{Alt}_NF)(x_1,\ldots,x_n)
 =\operatorname{sym}_N^{-1}
 \sum_{\sigma\in S_n}\operatorname{sgn}(\sigma)
 F(x_{\sigma(1)},\ldots,x_{\sigma(n)}).
\]
These maps commute with coefficient reduction in $N$.
\end{theorem}
\begin{proof}
Write $A=A_N$, $R=R_N$, and $A^e=A\tensor_R A^{\op}$.
In homological degree $n$, set
\[
 K_n=A\tensor_R\Lambda_R^n\flie_N\tensor_R A.
\]
Its augmentation is multiplication.
For $X=x_1\wedge\cdots\wedge x_n$, let $X_{\widehat i}$ and
$X_{\widehat i\widehat j}$ omit the indicated entries without reordering.
Define
\[
\begin{split}
 \partial(a\tensor X\tensor b)
 ={}&\sum_i(-1)^{i+1}
 \bigl(ax_i\tensor X_{\widehat i}\tensor b
       -a\tensor X_{\widehat i}\tensor x_i b\bigr)\\
 &+\hbarq\sum_{i<j}(-1)^{i+j}
 a\tensor[x_i,x_j]\wedge X_{\widehat i\widehat j}\tensor b.
\end{split}
\]
The normalized relative bimodule bar has terms
$\mathcal B_n=A\tensor_R\bar A^{\tensor_R n}\tensor_R A$.
Define
\[
 \iota_n(a\tensor X\tensor b)
 =\sum_{\sigma\in S_n}\operatorname{sgn}(\sigma)
       a[x_{\sigma(1)}|\cdots|x_{\sigma(n)}]b.
\]
The outer bar faces give the first sum in $\partial$.
Pair adjacent transpositions in the internal faces and use
$xy-yx=\hbarq[x,y]$ to obtain the second sum.
For two inputs this reads
\[
\begin{split}
 b\bigl(1[x|y]1-1[y|x]1\bigr)
 ={}&x[y]1-y[x]1+1[x]y-1[y]x\\
 &-\hbarq\,1[[x,y]]1.
\end{split}
\]
Thus $b\iota=\iota\partial$.
PBW makes $\flie_N$ a free summand of $\bar A$.
Since every factorial is invertible in $R$, $\iota$ is injective.
The identity $b^2=0$ then gives $\partial^2=0$.

Give $K_n$ filtration degree equal to its outer PBW degrees plus $n$.
The associated graded is the diagonal Koszul complex of
$\Sym_R(\flie_N)$.
Its differential uses differences of left and right polynomial generators.
A linear change of variables makes these independent polynomial variables.
They form a regular sequence even when $R$ has nilpotents.
For an augmented cycle, lift a primitive of its highest-degree symbol
and subtract its differential.
The filtration degree decreases, so this procedure terminates.
Hence $K_\bullet\to A$ is a free bimodule resolution.

The relative bar is also free over $A^e$, with its usual unit-insertion contraction.
The map $\iota$ lifts the identity of $A$.
To construct a reverse chain map, first lift the augmentation in
degree zero. At each higher degree its required boundary is a cycle;
exactness of $K_\bullet$ and projectivity of the bar term give a lift.
Each composite and the identity induce the same augmentation.
Their difference is therefore a boundary in degree zero, and the same
lifting argument, degree by degree, constructs a chain homotopy.
Applying $\operatorname{Hom}_{A^e}(-,A)$ gives a cochain homotopy equivalence.

The induced differential on
$\operatorname{Hom}_R(\Lambda_R^n\flie_N,A)$ is
\[
\begin{split}
 (\delta_N\phi)(x_1,\ldots,x_{n+1})
 ={}&\sum_i(-1)^{i+1}
 [x_i,\phi(x_1,\ldots,\widehat{x_i},\ldots,x_{n+1})]_A\\
 &+\hbarq\sum_{i<j}(-1)^{i+j}
 \phi([x_i,x_j],x_1,\ldots,\widehat{x_i},
                         \ldots,\widehat{x_j},\ldots,x_{n+1}).
\end{split}
\]
By \eqref{eq:vol2-scaled-pbw-action}, the commutator term also carries
one factor of $\hbarq$. Both sums therefore become $\hbarq$ times the
ordinary Chevalley differential.
Precomposition with $\iota$ gives $\operatorname{Alt}_N$.
All formulas commute with coefficient reduction.
\end{proof}

\section{The formal coefficient limit}
Put $R=\C\llbracket\hbarq\rrbracket$ and $A=\varprojlim_N A_N$.
PBW identifies its underlying module with $S_0\llbracket\hbarq\rrbracket$.
Each coefficient of $\hbarq$ is a polynomial.
We also write $U_\hbar(\flie)=A$ and
obtain the continuous module isomorphism
\[
 \operatorname{sym}_\hbar=\varprojlim_N\operatorname{sym}_N:
 S_0\llbracket\hbarq\rrbracket\xrightarrow{\sim}U_\hbar(\flie).
\]
Define
\[
 A^{\widehat\tensor_R n}=\varprojlim_N A_N^{\tensor_{R_N}n},
 \qquad
 C^\bullet_{\HH,\mathrm{cont}}(A/R,A)
 =\varprojlim_N C^\bullet_{\HH}(A_N/R_N,A_N).
\]
Equivalently, these are continuous $R$-multilinear normalized cochains
on the completed relative tensors.

\begin{corollary}[Formal relative Rees comparison]\label{cor:vol2-formal-relative-rees}
The compatible antisymmetrizations induce a quasi-isomorphism
\[
 C^\bullet_{\HH,\mathrm{cont}}(A/R,A)
 \longrightarrow
 \bigl(C^\bullet(\flie,S_0)\llbracket\hbarq\rrbracket,
                    \hbarq d_{\CE}\bigr).
\]
\end{corollary}
\begin{proof}
Coefficient reduction is surjective in each cochain degree.
For Hochschild cochains, lift outputs on the PBW input basis.
The Chevalley exterior input has finite rank, so its outputs also lift.
The comparison cones are acyclic and their transition maps are degreewise surjective.
The long exact sequence shows that every transition kernel is acyclic.
Choose a primitive of a compatible cycle at level zero.
Lift it one level and correct its differential by a primitive in the transition kernel.
Iteration gives compatible primitives at every order.
The inverse-limit cone is therefore acyclic.
\end{proof}

The coefficient ring is part of this result.
Absolute $\C$-linear cochains also detect variations of the coefficient parameter.
For $\flie=0$, the continuous absolute derivation $\partial_{\hbarq}$
on $R$ has no relative counterpart.
For nonabelian $\flie$, the associated graded for the $\hbarq$ filtration
has zero differential, whereas $d_{\CE}$ need not vanish.
The theorem asserts a quasi-isomorphism of complexes. A comparison of
cup products, braces and their higher coherences requires additional
maps and identities. The inverse
limit above changes only the coefficient order $N$, at a fixed
finite-dimensional Lie algebra. Passing to all Hamiltonian jets also
requires compatible transition maps for the chosen cochain complexes
and a proof that their limit computes the intended continuous carrier.
